\chapter{한계와 test 계획}

\section{남은 한계}
\begin{itemize}[leftmargin=*]
\item \textbf{현금흐름은 근사다.} 월별 납입 기록이 없어 받은 총액을 월 상환액으로 나눠 정상 납입 개월을 추정했다(2장).
\item \textbf{기대수익 점수는 1년 단위 근사다.} 조기상환과 부도 시점을 무시한다. 부도와 조기상환을 함께 모형화한 연구(Li 외, 2023, European Journal of Operational Research)가 개선 방향이다.
\item \textbf{안정성은 좋은 시절의 출렁임이다.} 2013--16년은 경기침체가 없었다. 위기 시기(2007--09)는 표본이 작고 A--D 등급만 볼 수 있었다(A가 가장 안정적이었으나 수준은 낮았다).
\item \textbf{소득은 자기 신고 값이고, 그 신호는 사라졌다.}
\item \textbf{개인이 살 수 있는 물량이 줄었다.} 조각 대출 비중이 2013년 72\%에서 2016년 27\%로 줄었다.
\item \textbf{같은 데이터를 오래 봤다}(8장).
\end{itemize}

\section{test 규칙 v1·\allowbreak{}v2 (test 전에 고정, git 태그)}
v1은 그대로 실행하고, v2는 30--34단계에서 나온 후보를 더한다(\ft{docs/test_protocol_v2.md}, 태그 \ft{test-protocol-v2}).

\begin{table}[H]\centering\small
\begin{tabularx}{\linewidth}{@{}p{45mm}L@{}}\toprule
가설 & 지지 기준(주 표본: test의 2013--16년 36개월 종결 대출)\\\midrule
H1 성과 귀속(v1) & 세부등급 분해에서 선별 > 0, 선별 > 배분, 3개 연도 이상 선별 > 0\\
H2 선별 가치(v1) & 기대수익 순위 선별 가치의 95\% 구간 하한 > 0\\
H3 국채 대비(v1) & 전부 승인이 발행월 90\% 이상에서 국채보다 높음\\
H4 회사채 대비(v1) & 선별하면 회사채 이하 발행월이 줄어듦\\
H5 등급 안 신호 & 기대수익 등급 중립 롱숏 90\% 구간 하한 > 0\\
H6 두 축 & 승인 30\%·\allowbreak{}50\% 모두 부도확률 순이 Sharpe 높고 기대수익 순이 수익 높음\\
H7 단순 규칙 & 계좌 개설 규칙이 부도확률 순을 Sharpe와 수익 둘 다에서 이김\\
H8 Sharpe 정의 & 세 정의 모두 부도확률 순 > 기대수익 순\\\bottomrule
\end{tabularx}
\end{table}

\paragraph{해석 규칙(미리 정함).}
\begin{enumerate}[leftmargin=*]
\item 만기 도달 비율이 0.9 미만인 연도는 따로 보고하고 근거로 쓰지 않는다(종결 편향, 3장).
\item 60개월은 36개월과 섞지 않는다.
\item H5가 실패하면 ``LC 가격이 효율화되었다''를 대안 결론으로 쓴다(연소득 신호가 사라진 것과 같은 방향).
\item 크기보다 방향과 구간으로 판정한다.
\item H7이 실패하면 ``train에서 보고 고른 규칙은 과대평가된다''를 그대로 보고한다.
\end{enumerate}

\section{test에서 무엇을 기대하나(판단, 계산 아님)}
\begin{table}[H]\centering\small
\begin{tabularx}{\linewidth}{@{}p{40mm}p{22mm}L@{}}\toprule
주장 & 유지 가능성 & 근거와 위험\\\midrule
등급 안 신호(H5, H2) & 높음 & 48개월 전부 양수, 해마다 커짐, 여러 조건에서 유지. 위험: test가 우리가 안 본 시기에 몰린 경우\\
두 축(H6, H8) & 중상 & 9개 승인비율·\allowbreak{}두 수익률 정의에서 같은 방향. 위험: 60개월 비중, 종결 편향\\
효과 크기 1\%p & 중간 & 같은 기간 안의 결과. 시기가 바뀌면 크기는 쉽게 변함\\
단순 규칙(H7) & 중간 & 사후 선택이라 과대평가 가능\\\bottomrule
\end{tabularx}
\end{table}

\section{발표에서의 한 문장}
\begin{keybox}[한 문장]
``돈은 같은 등급 안에 있다. 기대수익으로 고르면 가장 많이 벌고, 부도 위험이 낮은 쪽으로 고르면 과제의 Sharpe가 가장 높으며, 둘의 차이는 고위험 등급이 위험만큼 금리를 받지 못한 데서 온다. 우리는 이 결론을 미리 정한 규칙으로 test에서 한 번만 확인한다.''
\end{keybox}
